% Project Ghost — arXiv submission (v0.2.0, 2026-06-10)
% Author: Javier Menéndez Mateos <jfhelvetius@gmail.com>
% Compiles with pdflatex + bibtex (or biber if you switch \bibliographystyle).
%
%   pdflatex main
%   bibtex main
%   pdflatex main
%   pdflatex main
%
% This source is also the basis for RV 2026 / FMAS 2026 workshop
% submissions; swap \documentclass to the venue's required template
% (e.g. lncs.cls for Springer LNCS workshops).

\documentclass[11pt,a4paper]{article}

\usepackage[utf8]{inputenc}
\usepackage[T1]{fontenc}
\usepackage[english]{babel}
\usepackage{microtype}
\usepackage{amsmath, amssymb, amsthm}
\usepackage{booktabs}
\usepackage{tabularx}
\usepackage{array}
\usepackage{listings}
\usepackage{xcolor}
\usepackage{enumitem}
\usepackage[a4paper,margin=1in]{geometry}
\usepackage[colorlinks=true,linkcolor=blue!60!black,citecolor=blue!50!black,urlcolor=blue!60!black]{hyperref}

% --- Listings (Python / shell / JSON) -----------------------------------
\definecolor{codebg}{rgb}{0.97,0.97,0.97}
\lstset{
  basicstyle=\ttfamily\small,
  backgroundcolor=\color{codebg},
  frame=single,
  rulecolor=\color{black!20},
  numbers=none,
  breaklines=true,
  showstringspaces=false,
  columns=fullflexible,
  xleftmargin=0.5em,
  xrightmargin=0.5em,
}

% --- Theorem environments -----------------------------------------------
\newtheorem{theorem}{Theorem}
\newtheorem{corollary}[theorem]{Corollary}

% --- Title block --------------------------------------------------------
\title{Project Ghost: A Verifiable Safety-Property Surface\\
       for Autonomy Under Uncertainty}
\author{Javier Menéndez Mateos\\
        \texttt{jfhelvetius@gmail.com}\\
        \small Independent}
\date{v0.2.0 --- \today}

\begin{document}
\maketitle

\begin{abstract}
Safety claims in autonomy research are typically asserted in prose
and illustrated by simulation videos that the reader cannot re-run.
We describe \textbf{Project Ghost}, an open-source platform whose
primary contribution is \textbf{a citation pattern that lets a third
party verify any safety claim against the recorded run, byte-exact,
via a single shell command}: \texttt{pip install project-ghost==0.2.2},
then \texttt{ghost verify-properties --mcap <log>}. The pattern
composes seven existing ingredients --- ADRs, content-addressed
MCAP telemetry, pure-function verifiers, Hypothesis property tests,
CI gating, tagged releases, and OIDC-signed PyPI wheels --- into one
coherent reproducibility unit, with the underlying invariants
additionally checked by TLA+/TLC.

To exercise the pattern, we instantiate five safety properties for
the closed loop of a reference autonomy supervisor (BAUD-v1,
ERUR-v1, MD-v1, RLB-v1, FPB-v1). Each is stated in a binding ADR,
verified by a pure function over the MCAP, exercised by ${\sim}50$
property tests within a 1687-test suite, witnessed inline in every
reference smoke, and self-enforced on every push by CI. Three TLA+
specifications jointly cover the property set with 11 invariants
verified by TLC on a bounded abstract model, including a partition
theorem $\mathrm{BAUD} \oplus \mathrm{ERUR}$ and a closed-form
recovery latency bound $L \le \mathrm{peak} + W - 1$
(Theorem~\ref{thm:rlb}), shown tight by a witness trace. We
acknowledge the latency bound is elementary in hindsight; we
present it as supporting evidence for the broader citation pattern
rather than as a stand-alone theoretical contribution.

Empirical evaluation on a violation matrix of six injected bug
categories, three structurally distinct calibration policies, three
shape-realistic drift profiles, and a head-to-head benchmark against
RTAMT establishes that the verifier is policy-agnostic, runs in
21--406\,ms (linear in trace length), and produces byte-identical
MCAPs and canonicalised property-report JSON across Linux and
Windows CI runners. The full artifact is re-runnable from
\texttt{pip install project-ghost==0.2.2}.
\end{abstract}

\medskip
\noindent\textbf{Keywords:} safety citation patterns, reproducible
safety verification, runtime verification, content-addressed
telemetry, calibrated confidence, TLA+/TLC, MCAP.

\medskip
\noindent\textbf{Artifact availability:}
\begin{itemize}[noitemsep,topsep=0pt]
  \item Source: \url{https://github.com/JFHelvetius/ghost}
  \item Distribution: \url{https://pypi.org/project/project-ghost/}
  \item Documentation: \url{https://JFHelvetius.github.io/ghost/}
  \item License: Apache-2.0
\end{itemize}

% =======================================================================
\section{Introduction}
\label{sec:intro}

Safety claims in robotics are routinely supported by hand-written
prose in design documents and by simulation videos that the reader
cannot re-run. The literature on uncertainty in autonomy is rich ---
Bayesian filters~\cite{thrun2005probabilistic}, calibration of
probabilistic predictions~\cite{vovk2005algorithmic}, epistemic
versus aleatoric uncertainty~\cite{kendall2017uncertainties}, fault
detection and isolation (FDI)~\cite{isermann2006fdi}, runtime safety
supervisors~\cite{bartocci2018lectures} --- but the gap between
\emph{the theory exists} and \emph{this specific run, on this
specific code, satisfies the property} is rarely operationally
closed. A third party who wants to verify a safety claim against a
recorded run typically cannot: there is no shell command, no
content-addressed log, no pure-function verifier they can re-run on
their own machine.

We describe Project Ghost, an opinionated reference platform built
around exactly this gap. Ghost is sim-first, written in Python, and
ships as a \texttt{pip}-installable package with a CLI subcommand
(\texttt{ghost verify-properties}) that takes a captured MCAP log
and returns a byte-exact verdict on five formal safety properties.
Each property is stated as a binding ADR; verified by a pure function
over the log; exercised by Hypothesis-based property tests;
witnessed inline in every reference closed-loop smoke; and
self-enforced on every push by CI. Two of the properties (BAUD-v1 and
ERUR-v1) are additionally \textbf{mechanically verified} by TLA+/TLC
over the abstract state space of the reference policy pair, along
with the partition theorem that the two properties together cover
the full conditional behaviour space. The tight recovery-latency
bound RLB-v1 is mechanically verified by a separate TLA+ spec
mirroring the verifier algorithm.

\subsection{Contributions}
\label{sec:contributions}

We order our contributions \textbf{from the most load-bearing (an
engineering pattern) to the most narrow (a useful supporting
bound)}:

\begin{itemize}
  \item \textbf{C1 --- A safety citation pattern.} A composition of
        seven existing ingredients --- ADR + content-addressed MCAP +
        pure-function verifier + Hypothesis property test + CI gate +
        tagged release + OIDC-signed PyPI wheel --- assembled as a
        single reproducibility unit so a third party can verify any
        cited safety claim against the captured run via one shell
        command, without trusting the producer. This is the pattern
        we believe is genuinely differentiating; the rest of the
        paper is the evidence that it works in practice.
  \item \textbf{C2 --- A reproducibility primitive with demonstrated
        detection capacity.} A one-line CLI verifier
        \texttt{ghost verify-properties} over content-addressed MCAP
        logs, distributed via PyPI with OIDC trusted publishing.
        Bug-detection is demonstrated systematically in
        \S\ref{sec:violation} via a six-category violation matrix
        (injected calibrator, decision, actuation, and threshold
        bugs; all six detected). The verifier produces deterministic
        JSON output across Linux and Windows runners and remains
        policy-agnostic across three structurally distinct
        calibration policies (\S\ref{sec:multi-policy}).
  \item \textbf{C3 --- A property set with mechanically-checked
        semantics for a reference autonomy supervisor.} Five
        properties (BAUD-v1, ERUR-v1, MD-v1, RLB-v1, FPB-v1)
        instantiating the citation pattern on a representative closed
        loop. Three TLA+ specifications cover the property set with
        11 invariants verified by TLC on a bounded abstract model in
        CI, including a partition theorem
        $\mathrm{BAUD} \oplus \mathrm{ERUR}$. The properties
        themselves are deliberately simple; the contribution is the
        end-to-end mechanisation, not the formulation.
  \item \textbf{C4 --- A closed-form tight recovery latency bound}
        for sliding-window count-of-$K$-in-$W$ filters
        (Theorem~\ref{thm:rlb}): $L \le \mathrm{peak} + W - 1$,
        attained with equality by a witness trace
        ($L = 38 = 7 + 32 - 1$). The bound follows directly from the
        sliding-window mechanism and is elementary in hindsight; we
        did not locate it stated as a closed form in our literature
        review (\S\ref{sec:related}) but we make no broader novelty
        claim and treat it as supporting evidence for C3 rather than
        a stand-alone theoretical result.
\end{itemize}

C1 and C2 are the contributions we expect to age best. C3
instantiates them on a representative supervisor; C4 supplies a
useful bound that the spec \texttt{Rlb.tla} verifies mechanically.
We position the work as a \textbf{systems / tools paper, not a
theory paper}.

\subsection{What this paper is and is not}

This is an engineering and infrastructure paper, not a theory paper.
The filtering, calibration, and FDI ingredients Ghost rests on are
well established. Theorem~\ref{thm:rlb}'s mathematics is elementary
in hindsight, and we present it as evidence that the citation
pattern can carry a sharp result, not as a stand-alone theoretical
contribution. The partition theorem of \S\ref{sec:mechanical} is
novel \emph{in the form we mechanised it} --- a TLA+
\texttt{INV\_PARTITION} over the reference closed loop --- but the
underlying observation that ``drift fired'' and ``drift clean and
KNOWN'' partition the per-cycle condition space is structurally
simple. We deliberately resist overclaiming on either result. The
contribution we are most willing to defend is \textbf{the engineered
combination} (C1) and its \textbf{operational demonstration} (C2):
that a third party can issue one shell command and obtain a
byte-exact verdict against a binding ADR.

% =======================================================================
\section{Background and related work}
\label{sec:related}

\subsection{Underlying ingredients}

Project Ghost builds on standard practice in robotics and control:
Bayesian and particle filtering~\cite{thrun2005probabilistic};
calibration of probabilistic
predictions~\cite{vovk2005algorithmic}; epistemic vs.\ aleatoric
uncertainty~\cite{kendall2017uncertainties}; fault detection and
isolation~\cite{isermann2006fdi}; runtime
verification~\cite{bartocci2018lectures}; TLA+ and TLC for
explicit-state model checking~\cite{lamport2002specifying}; MCAP for
content-addressable robotics telemetry~\cite{foxglove2022mcap}.

\subsection{Closest tooling prior work}

\begin{itemize}
  \item \textbf{RTAMT}~\cite{nikovic2020rtamt,nikovic2023rtamt}: STL
        monitors over CPS logs with online/offline algorithms and a
        Python API. STL-defined properties, no mechanically verified
        proof layer, no content-addressed reproducibility chain.
  \item
        \textbf{MoonLight}~\cite{bartocci2020moonlight,bartocci2023moonlight}:
        STREL (spatio-temporal logic) monitor in Java with a CLI.
        Spatial focus; no formal verification of monitor semantics.
  \item
        \textbf{ROSMonitoring}~\cite{ferrando2020rosmonitoring} and
        \textbf{ROSRV}~\cite{huang2014rosrv}: live ROS-middleware
        monitors. Both online; neither performs post-hoc log
        verification with a one-line CLI.
  \item \textbf{Safe RL via
        shielding}~\cite{jansen2020shielding,jansen2024shielding}:
        runtime enforcement of safety via action filters. Online,
        action-blocking; Ghost is offline and log-verifying.
  \item \textbf{Control Barrier Functions}~\cite{mitll2023cbf}:
        controller synthesis for continuous safety constraints.
        Complementary, not competing.
  \item \textbf{Conformal prediction for robot
        safety}~\cite{chakraborty2024conformal}: forward-looking
        distribution-free uncertainty bounds for gating actions.
        Predictive; Ghost is retrospective.
  \item \textbf{Supervisory control of timed
        automata}~\cite{flordal2022timedsupervisory}: synthesises
        timed supervisors. Constructs new supervisors; Ghost verifies
        existing traces. Prior timed-automata recovery results do
        not give the closed-form bound of Theorem~\ref{thm:rlb}.
  \item \textbf{Surveys of formal verification for
        autonomy}~\cite{rizaldi2020survey}: catalogue
        Coq/Lean/Isabelle/Alloy work. Note the absence of
        mechanically-verified TLA+ specs for autonomy supervisors
        specifically.
\end{itemize}

\subsection{Comparison matrix}
\label{sec:comparison}

\begin{table}[h]
\centering
\small
\setlength{\tabcolsep}{4pt}
\begin{tabularx}{\linewidth}{l X X X X X}
\toprule
\textbf{Dimension} & \textbf{Ghost} & \textbf{RTAMT} & \textbf{MoonLight} & \textbf{Shielding} & \textbf{Timed Aut.~SC} \\
\midrule
Verification mode             & Post-hoc log              & On/offline      & On/offline      & Online enforce & Offline synth.\\
Distribution                  & PyPI + OIDC               & Source          & Source          & Framework       & Synth. tool   \\
Content-addressed input       & \textbf{Yes} (SHA-256)    & No              & No              & N/A             & No            \\
One-line CLI verifier         & \textbf{Yes}              & No              & No              & No              & No            \\
Property nature               & Behavioural + latency     & STL             & STREL           & Invariants      & Discrete/timed\\
Mechanical proof              & \textbf{TLA+/TLC}         & None            & None            & Informal        & Timed-aut.    \\
Multi-property output         & \textbf{5 reports/run}    & 1/spec          & 1/spec          & Modular         & 1/synth.      \\
Partition theorem             & \textbf{BAUD$\oplus$ERUR} & N/A             & N/A             & N/A             & N/A           \\
Closed-form recovery bound    & \textbf{$L\le \mathrm{peak}+W-1$} & N/A     & N/A             & N/A             & None          \\
Bug-detection demo            & \textbf{Yes (\S\ref{sec:violation})} & N/A  & N/A             & N/A             & N/A           \\
\bottomrule
\end{tabularx}
\caption{Comparison of Project Ghost with the closest runtime
         verification and supervisory-control tooling. To our
         knowledge, no prior tool ships a content-addressed,
         pure-function safety-property verifier via \texttt{pip
         install} with mechanically verified underlying invariants.}
\label{tab:comparison}
\end{table}

% =======================================================================
\section{The property set}
\label{sec:properties}

The five properties are stated in binding ADRs (immutable once
accepted) and verified by pure functions in
\texttt{src/project\_ghost/properties/}. Each verifier returns a
typed report with \texttt{holds: bool}, structured per-cycle
metadata, and the MCAP's SHA-256.

\begin{itemize}
  \item \textbf{BAUD-v1 (ADR-0031): Bounded Action Under Drift.}
        Over a sliding window of size $W=32$, at least $M=4$
        calibration outcomes observed and at least $K=2$ in the
        dirty band. Postcondition: adjusted level strictly lower
        than raw in the lattice; emitted decision is not PROCEED;
        emitted actuator command, if any, belongs to a closed
        safe-reason set
        $S_{\mathrm{BAUD}} = \{\texttt{attitude\_hold\_hold}, \texttt{kill\_zero\_throttle}\}$.
  \item \textbf{ERUR-v1 (ADR-0032): Eventual Reactivation Under
        Recovery.} Drift absent and raw belief is KNOWN.
        Postcondition: adjusted level is KNOWN and decision is
        PROCEED. Together with BAUD, ERUR forms the
        \textbf{partition theorem}.
  \item \textbf{MD-v1 (ADR-0033): Monotonic Degradation.} For every
        cycle, $\mathrm{adjusted} \preceq \mathrm{raw}$ in the
        confidence lattice. The calibration policy never \emph{invents}
        confidence.
  \item \textbf{RLB-v1 (ADR-0034): Recovery Latency Bound.} See
        Theorem~\ref{thm:rlb}.
  \item \textbf{FPB-v1 (ADR-0035): False Positive Bound observer.}
        Empirical BAUD fire rate over the run, exposed as a
        structured metric for regression gating.
\end{itemize}

% =======================================================================
\section{Verifier architecture}
\label{sec:verifier}

Each property has a verifier
\texttt{src/project\_ghost/properties/verify\_<id>.py} that opens the
MCAP read-only, walks the channels of interest in cycle order,
computes the precondition and postcondition per cycle from the
stored messages alone, and returns a typed report. No replay, no
simulation. The report dataclasses
(\texttt{BAUDVerificationReport}, \texttt{ERURViolation}, etc.) are
public API; their JSON serialisation is the wire format the CLI
emits with \texttt{--json}.

\begin{lstlisting}[language=bash]
$ pip install project-ghost==0.2.2
$ python -m project_ghost.examples.closed_loop_smoke
$ ghost verify-properties --mcap closed_loop_smoke.mcap
BAUD-v1: HOLDS  (M=4, K=2, 6/10 cycles evaluated)
ERUR-v1: HOLDS  (M=4, K=2, 4/10 cycles evaluated)
MD-v1:   HOLDS  (10/10 cycles evaluated)
RLB-v1:  HOLDS  (W=32, 0/10 cycles evaluated)
FPB-v1:  HOLDS  (fire_fraction=0.60, 6/10 cycles evaluated)
$ echo $?
0
\end{lstlisting}

Exit code conventions: \texttt{0} iff every property holds, \texttt{1}
if any property violates or the verifier crashes, \texttt{2} for
argument errors. \texttt{--json} emits a deterministic JSON object
suitable for CI consumption.

\paragraph{CI as continuous guarantee.}
\texttt{.github/workflows/ci.yml} includes a
\texttt{verify-properties} job that runs the smoke and the verifier
on every push; a \texttt{tla-plus} job that runs TLC on the
mechanical specs (\S\ref{sec:mechanical}) on every push; a
\texttt{determinism-cross-machine} job that asserts byte-equality of
the smoke MCAP and the verifier JSON output across Linux and Windows
runners; and a \texttt{release} workflow that builds the wheel,
installs it in a fresh venv, re-runs the verifier against the
bundled smoke MCAP from the \emph{installed} wheel, and publishes
to PyPI via OIDC trusted publishing only if everything is green.

% =======================================================================
\section{Mechanical verification}
\label{sec:mechanical}

\subsection{Why TLA+}
Property-based testing with Hypothesis (200+ examples per property)
provides strong evidence at production scale, but proves the property
holds on the inputs the generator sampled, not on all inputs.
Mechanical verification over a finite abstract model closes the
remaining gap. We pick TLA+ with TLC over theorem
proving~\cite{lamport2002specifying} on a cost/benefit argument: TLC
is exhaustive over a finite state space in seconds, where a Lean
proof would be weeks.

\subsection{Three specifications, jointly covering 5/5 properties}
\textbf{\texttt{BaudErur.tla}} (\texttt{docs/proofs/BaudErur.tla},
bounds $M{=}2, K{=}1, W{=}3$) models the closed loop as a state
machine with one transition per cycle. It checks five invariants:
\texttt{INV\_BAUD}, \texttt{INV\_ERUR}, \texttt{INV\_PARTITION}
(contribution C2), \texttt{INV\_NO\_INVENTED\_CONFIDENCE} (MD-v1),
and \texttt{INV\_HISTORY\_BOUND}.

\textbf{\texttt{Rlb.tla}} (\texttt{docs/proofs/Rlb.tla}, bounds
$W{=}4, \mathrm{MAX\_DRIFT}{=}4$) mirrors the verifier algorithm of
\texttt{src/project\_ghost/properties/rlb.py} under the
consecutive-drift hypothesis and checks three invariants:
\texttt{INV\_RLB} (the mechanical witness of Theorem~\ref{thm:rlb}),
\texttt{INV\_PEAK\_BOUNDED}, and \texttt{INV\_WINDOW\_BOUND}.

\textbf{\texttt{Fpb.tla}} (\texttt{docs/proofs/Fpb.tla}, bounds
$\mathrm{MAX\_CYCLES}{=}8$, $\mathrm{MAX\_FIRE\_NUMER}{=}\mathrm{BOUND\_DENOM}{=}1$)
models the FPB-v1 counter automaton and checks three invariants on
the verifier's structural well-formedness:
\texttt{INV\_FPB\_RATIO\_BOUNDED}
($\mathrm{cycles\_fires} \le \mathrm{cycles\_total}$),
\texttt{INV\_FPB\_FIRE\_IMPLIES\_TOTAL}, and
\texttt{INV\_FPB\_OBSERVATIONAL\_DEFAULT} (the bound trivially holds
under the default $\mathrm{max\_fire\_fraction} = 1.0$). FPB is
purely observational; a statistical FPB-v2 with Monte Carlo bounds
is future work.

Together the three specs constitute \textbf{5/5 properties with at
least a structural TLC invariant in CI}.

\subsection{Scope}
The mechanical proofs do not claim that the Python implementation
faithfully mirrors the TLA+ model (the bridge is by inspection),
nor that the bounded constants prove the unbounded case. Behaviour
at production-scale constants is covered by the property tests.

% =======================================================================
\section{A closed-form recovery latency bound}
\label{sec:theorem}

\subsection{Setting}
Let $(o_t)_{t\ge 1}$ be the stream of per-cycle prediction
outcomes, classified into a binary partition
$\mathrm{dirty}(o) \in \{0,1\}$ where $\mathrm{dirty}(o) = 1$ when
the Mahalanobis verdict is at or above the BAUD threshold. Let
$H_t$ denote the sliding window of the last $W$ outcomes available
at cycle $t$:
\begin{equation}
H_t = (o_{\max(1, t-W+1)}, \ldots, o_t), \qquad |H_t| \le W.
\end{equation}
The reference calibrator (\texttt{MahalanobisDowngradePolicy}$(M,K)$)
downgrades the adjusted self-assessment level on any cycle where
\begin{equation}\label{eq:cond1}
|H_t| \ge M \quad \text{and} \quad \sum_{o \in H_t} \mathrm{dirty}(o) \ge K.
\end{equation}

\subsection{Theorem 1 and proof}

\begin{theorem}[Tight Recovery Latency Bound; transient regime]\label{thm:rlb}
Let the outcome stream contain a transient drift interval of
$N \le W$ consecutive dirty outcomes followed by clean outcomes,
under a sliding window of size $W$. Define
\begin{itemize}[noitemsep,topsep=2pt]
\item $\mathrm{peak} = \min(N, W) = N$, the maximum dirty count
      observed in the window during the drift run;
\item $L$, the dirty-run length: the number of consecutive cycles
      where the window contains at least one dirty outcome.
\end{itemize}
Then $L = \mathrm{peak} + W - 1$. Equivalently, the bound
$L \le \mathrm{peak} + W - 1$ is attained with equality.
\end{theorem}

\begin{proof}
Trace the window state cycle by cycle, noting the sliding-window
invariant: at cycle $t$, the window contains the last
$\min(t, W)$ outcomes.
\begin{itemize}
  \item \textbf{Accumulation} (cycles $1..N$). Each cycle adds one
        dirty outcome; the window is not yet full ($N \le W$), so
        no expulsion. The dirty count rises from $1$ to
        $N = \mathrm{peak}$. All $N$ cycles are dirty (count $\ge 1$).
  \item \textbf{Saturation} (cycles $N+1..W$). Each cycle adds one
        clean outcome; the window is still not full, so no
        expulsion. Dirty count stays at $\mathrm{peak}$. All
        $W - N$ cycles are dirty.
  \item \textbf{Flush} (cycles $W+1..W+\mathrm{peak}-1$). The
        window is now full; each new clean outcome expels the
        oldest entry. By construction, the oldest entries are the
        dirty outcomes that arrived first. Dirty count decreases
        by 1 per cycle, from $\mathrm{peak}$ to $1$. All
        $\mathrm{peak} - 1$ cycles are dirty.
  \item \textbf{Recovery} (cycle $W+\mathrm{peak}$). The last
        dirty outcome is expelled. Dirty count is $0$. This cycle
        is clean.
\end{itemize}
Summing dirty cycles: $N + (W - N) + (\mathrm{peak} - 1) = W + \mathrm{peak} - 1$.
Since $\mathrm{peak} = N$ in the transient regime,
$L = \mathrm{peak} + W - 1$. \qedhere
\end{proof}

\begin{corollary}[Operational regime]\label{cor:regime}
When $N > W$, the drift outlasts the window; $\mathrm{peak} = W$
and $L = N + W - 1$. The bound $\mathrm{peak} + W - 1 = 2W - 1$ is
exceeded whenever $N > W$. The bound therefore operationally
characterises the \emph{transient} regime; in the sustained-drift
regime, no recovery transition occurs \emph{during the drift} and
the verifier records the property vacuously.
\end{corollary}

\begin{corollary}[Structural sanity]\label{cor:structural}
A trace where $L > \mathrm{peak} + W - 1$ on a recovery transition
is impossible under a correctly implemented sliding window of
size $W$. The verifier's \texttt{RLBViolation} therefore also
serves as a structural-integrity check on the window
implementation.
\end{corollary}

\subsection{Operational tightness check}
The drift-then-recovery smoke
\texttt{closed\_loop\_smoke\_with\_recovery.py} is engineered to
exhibit Theorem~\ref{thm:rlb} at the production constants
($N = \mathrm{peak} = 7$, $W = 32$):
$L_{\mathrm{observed}} = 38 = 7 + 32 - 1$. The integration test
\texttt{tests/integration/test\_closed\_loop\_smoke\_with\_recovery.py}
asserts the recovery transition fires at exactly cycle 39 and
nowhere earlier or later.

% =======================================================================
\section{Reproducibility surface}
\label{sec:reproducibility}

The headline claim is that a third party can verify the property
set against a captured run \emph{without trusting the producer}.
The reproducibility surface has five layers: (1) content-addressed
MCAP with SHA-256 in every property report; (2) deterministic
pipeline (ADR-0030 Replay Verification v1); (3) pure-function
verifiers with no I/O beyond reading the MCAP; (4) approximately 50
Hypothesis-based property tests; (5) TLA+ continuous self-check on
every push.

A reader who wishes to cite a Project Ghost safety claim can write:
\begin{quote}\small
Project Ghost v0.2.0 satisfies BAUD-v1 on the bundled reference
smoke MCAP \texttt{SHA-256:<hash>}, as verified by
\texttt{ghost verify-properties --mcap closed\_loop\_smoke.mcap}
from \texttt{pip install project-ghost==0.2.2}, and additionally
satisfies \texttt{INV\_BAUD}, \texttt{INV\_ERUR},
\texttt{INV\_PARTITION} over the abstract model
\texttt{BaudErur.tla} at bounds $M{=}2, K{=}1, W{=}3$, and
\texttt{INV\_RLB} (Theorem~\ref{thm:rlb}) over
\texttt{Rlb.tla} at $W{=}4$.
\end{quote}
This is contribution C4 in action.

% =======================================================================
\section{Evaluation}
\label{sec:evaluation}

\subsection{Tests, CI, and mechanical verification}
At v0.2.0, the test suite contains 1665 tests passing (ruff +
mypy strict + deptry clean), of which approximately 50 are
dedicated property tests. The CI matrix runs on
\texttt{ubuntu-latest} and \texttt{windows-latest} with Python
3.11 and 3.12, plus a \texttt{tla-plus} job that runs TLC on both
mechanical specs on every push and uploads
\texttt{tlc\_output\_*.log} as a build artifact.

\subsection{Bug-detection capability (Violation Matrix)}
\label{sec:violation}

A standalone smoke \texttt{closed\_loop\_smoke\_violated.py}
demonstrates the verifier detects \emph{a} bug. To demonstrate that
detection is \textbf{systematic, not anecdotal}, we extend this to
a violation matrix of six bug categories, one mini-smoke per
category, each engineered to break exactly one component of the
closed loop.

\begin{table}[h]
\centering
\small
\begin{tabular}{lll c}
\toprule
\textbf{Category} & \textbf{Buggy component} & \textbf{Property expected} & \textbf{Detected?} \\
\midrule
\texttt{calibrator\_no\_downgrade}        & calibration policy           & BAUD-v1               & YES \\
\texttt{calibrator\_invents\_confidence}  & calibration policy           & MD-v1 (and BAUD-v1)   & YES \\
\texttt{decision\_proceeds\_anyway}       & decision policy              & BAUD-v1               & YES \\
\texttt{decision\_never\_proceeds}        & decision policy              & ERUR-v1               & YES \\
\texttt{actuation\_non\_safe\_reason}     & actuation policy             & BAUD-v1               & YES \\
\texttt{fpb\_threshold\_exceeded}         & verifier \texttt{max\_fire\_fraction} & FPB-v1     & YES \\
\bottomrule
\end{tabular}
\caption{Violation matrix: six bug categories, each engineered to
break exactly one component of the closed loop, all six detected
by the unmodified verifier. The matrix covers BAUD-v1, ERUR-v1,
MD-v1, and FPB-v1; RLB-v1 is structural to the sliding-window
implementation and is exercised by the drift-then-recovery smoke at
the bound (\S\ref{sec:theorem}). Reproducible via
\texttt{python -m project\_ghost.examples.violation\_matrix};
the script's exit code is 1 iff any false negative occurs.}
\label{tab:violation-matrix}
\end{table}

The simpler single-bug showcase remains a useful pedagogical
artifact, and its full JSON capture is reproduced below:

\begin{lstlisting}[language=bash]
$ ghost verify-properties --mcap closed_loop_smoke_violated.mcap --json
{
  "all_properties_hold": false,
  "properties": {
    "BAUD-v1": { "holds": false, "violation_count": 12,
                 "cycles_precondition_held": 6, "cycles_total": 10 },
    "ERUR-v1": { "holds": true,  ... },
    "MD-v1":   { "holds": true,  ... },
    "RLB-v1":  { "holds": true,  ... },
    "FPB-v1":  { "holds": true,  ... }
  }
}
$ echo $?
1
\end{lstlisting}

The verifier detects 12 individual postcondition violations across
6 cycles where BAUD's precondition fires (6 cycles $\times$ 2
postconditions). Exit code 1. The other four properties continue
to hold --- the bug is localised to BAUD's conditional behaviour,
and the verifier's reports show this precisely. This demonstrates
detection capacity for contribution C3 and differential
failure-mode visibility across the property set.

\subsection{Parametric policy evaluation}
\label{sec:parametric}

To demonstrate stability under variation of the reference
calibrator parameters, we ran the closed-loop smoke under three
$(M, K)$ pairs and three trace lengths
$n \in \{10, 50, 200\}$. All 9 combinations pass all 5 properties.
Verifier runtime is linear in trace length and policy-insensitive.
Reproducible via
\texttt{docs/paper/scripts/measure\_metrics.py}.

\begin{table}[h]
\centering
\small
\begin{tabular}{lrrrrrc}
\toprule
\textbf{Policy $(M,K)$} & $n$ & \textbf{MCAP (B)} & \textbf{Smoke (ms)} & \textbf{Verifier total (ms)} & \textbf{BAUD fire frac} & \textbf{All HOLDS} \\
\midrule
$(4,2)$ reference   & 10  & 6\,552  &  15.3 & 20.8  & 0.60 & \checkmark\\
$(4,2)$ reference   & 50  & 18\,481 &  24.3 & 99.7  & 0.92 & \checkmark\\
$(4,2)$ reference   & 200 & 64\,699 & 100.5 & 405.8 & 0.98 & \checkmark\\
$(3,1)$ sensitive   & 10  & 6\,542  &   5.6 & 20.6  & 0.70 & \checkmark\\
$(3,1)$ sensitive   & 50  & 18\,475 &  34.0 & 100.2 & 0.94 & \checkmark\\
$(3,1)$ sensitive   & 200 & 64\,717 & 102.7 & 399.1 & 0.99 & \checkmark\\
$(5,3)$ lax         & 10  & 6\,553  &   5.6 & 20.6  & 0.50 & \checkmark\\
$(5,3)$ lax         & 50  & 18\,484 &  23.4 & 100.7 & 0.90 & \checkmark\\
$(5,3)$ lax         & 200 & 64\,700 & 104.4 & 406.2 & 0.98 & \checkmark\\
\bottomrule
\end{tabular}
\caption{Parametric evaluation across three calibrator
         parameterisations and three trace lengths. Verifier
         timing is the mean of 5 replicates with one warm-up.
         All 9 combinations report all 5 properties HOLDS.}
\label{tab:parametric}
\end{table}

\subsection{Policy-agnostic verifier, policy-specific preconditions}
\label{sec:multi-policy}

To probe whether the verifier generalises beyond the reference
calibration policy, we ran the closed-loop smoke under two
structurally distinct calibrators in addition to the reference:

\begin{itemize}[noitemsep,topsep=2pt]
  \item \texttt{MahalanobisDowngradePolicy}$(M{=}4, K{=}2)$ ---
        reference; Theorem~\ref{thm:rlb} applies.
  \item \texttt{EWMADowngradePolicy}$(\alpha{=}0.5, \min{=}3, \mathrm{thr}{=}0.3)$
        --- exponentially-weighted moving average over the dirty
        indicator. Theorem~\ref{thm:rlb} does not apply.
  \item \texttt{PerAxisHysteresisDowngradePolicy}$(\mathrm{up}{=}3.0)$
        --- per-axis Mahalanobis distance with hysteresis.
\end{itemize}

The verifier was executed unchanged on all three resulting MCAPs.
Per-property results (verifier parameterised with the
\emph{reference's} $M{=}4, K{=}2$):

\begin{table}[h]
\centering
\small
\begin{tabular}{lccccc}
\toprule
\textbf{Policy} & \textbf{BAUD} & \textbf{ERUR} & \textbf{MD} & \textbf{RLB} & \textbf{FPB} \\
\midrule
\texttt{MahalanobisDowngradePolicy}$(M{=}4, K{=}2)$ ref. & OK & OK & OK & OK & OK \\
\texttt{EWMADowngradePolicy}$(\alpha{=}0.5)$            & OK & \textbf{VIOL} & OK & OK & OK \\
\texttt{PerAxisHysteresisDowngradePolicy}              & OK & \textbf{VIOL} & OK & OK & OK \\
\bottomrule
\end{tabular}
\caption{Policy-comparison matrix. The verifier reads each MCAP
without modification; ERUR-v1 violations on the alternative
policies are informative, not bugs (see discussion).}
\label{tab:policy-comparison}
\end{table}

\textbf{The ERUR violations are informative, not bugs.} Three
observations: (1) the verifier is genuinely policy-agnostic ---
all three policies share the same Protocol and MCAP schema;
(2) MD-v1 holds on all three policies by independent obligation
(monotonicity is enforced per-policy); (3) BAUD-v1 and ERUR-v1
are parameterised on the reference policy's $(M, K)$, so running
the verifier with reference parameters against EWMA / PerAxis
traces asks ``did this trace behave like a reference $(M{=}4,
K{=}2)$ policy would have?'' --- the alternative policies do not
always agree, and ERUR therefore sometimes violates.

The contribution that flows from this observation: the verifier
remains useful as a third-party safety check for any calibrator,
\emph{but the property parameters must match the operator's policy
contract}, not blindly inherit the reference's. A future ADR-0037
could state ERUR abstractly over \texttt{policy.precondition(history)};
v0.2.x keeps the property concrete and visible.

\subsection{Realistic-shape scenarios}
\label{sec:realistic}

To demonstrate the verifier handles drift profiles beyond the
sustained-drift / single-recovery patterns, we run three additional
scenarios whose ground-truth functions are shaped after failure
modes documented in the VIO/SLAM evaluation literature.

\begin{itemize}[noitemsep,topsep=2pt]
  \item \texttt{gps\_denial}: 6 drift cycles + 44 recovery
        (transient regime; Theorem~\ref{thm:rlb} applies).
  \item \texttt{slow\_biased\_drift}: 50 cycles of low-magnitude
        chronic drift; FPB fire fraction climbs to 0.92.
  \item \texttt{cascading\_failure}: 30 cycles in three phases —
        stationary, position-drift, position+yaw-drift.
\end{itemize}

\begin{table}[h]
\centering
\small
\begin{tabular}{lrcccccr}
\toprule
\textbf{Scenario} & \textbf{Cycles} & \textbf{BAUD} & \textbf{ERUR} & \textbf{MD} & \textbf{RLB} & \textbf{FPB} & \textbf{fire\_frac} \\
\midrule
\texttt{gps\_denial}         & 50 & OK & OK & OK & OK & OK & 0.64 \\
\texttt{slow\_biased\_drift} & 50 & OK & OK & OK & OK & OK & 0.92 \\
\texttt{cascading\_failure}  & 30 & OK & OK & OK & OK & OK & 0.60 \\
\bottomrule
\end{tabular}
\caption{Realistic-shape scenarios: all five properties hold under
the reference policy. Reproducible via
\texttt{python -m project\_ghost.examples.realistic\_scenarios}.}
\label{tab:realistic}
\end{table}

\textbf{Honest scope: shape-realistic, not data-real.} The
scenarios are deterministic synthetic profiles informed by the
literature, not extracted from real flight telemetry. A genuine
flight-data evaluation requires an adapter from PX4 ULog or
ROSBag to the Ghost pipeline inputs; the roadmap is
\texttt{docs/paper/venues/dataset\_integration.md}. The scenarios
here establish that the verifier and property set generalise to
non-trivial failure shapes; full flight-data validation is future
work.

\subsection{Quantitative comparison against RTAMT}
\label{sec:rtamt-bench}

To put the comparison matrix of \S\ref{sec:comparison} on
quantitative ground, we benchmarked Ghost's \texttt{verify\_baud}
against an STL-based safety check in RTAMT~\cite{nikovic2020rtamt}
(version 0.3.5 from PyPI). The benchmark generates a fresh 50-cycle
smoke MCAP, extracts time-aligned signals
(\texttt{error\_magnitude}, \texttt{proceed\_indicator}), and times
both verifiers over 5 replicates with a warm-up cycle.

The STL formula approximating BAUD-v1 is
$G_{[0,0.4]}(\mathrm{error} > 3) \rightarrow (\mathrm{proceed} < 0.5)$:
``whenever the prediction error has been above 3-sigma for the past
0.4-second window ($\sim$4 cycles at 10\,Hz), the agent must not be
issuing PROCEED''. This is the closest expressible approximation;
STL cannot natively count $K$-of-$M$-in-$W$ occurrences as a single
formula.

\begin{table}[h]
\centering
\small
\begin{tabular}{lp{4.2cm}cr l}
\toprule
\textbf{Tool} & \textbf{Property language} & \textbf{Verdict} & \textbf{Time (ms)} & \textbf{Setup} \\
\midrule
Ghost \texttt{verify\_baud} (BAUD-v1 exact) & Python predicate over MCAP schema & \textbf{HOLDS} & 23.44 & 1-line CLI \\
RTAMT (STL dense-time)                     & $G_{[0,0.4]}(e>3 \to p<0.5)$       & VIOLATED         &  0.15 & 12-line glue \\
\bottomrule
\end{tabular}
\caption{Quantitative benchmark on the reference 50-cycle smoke
MCAP. Wall-clock times are means of 5 replicates with a warm-up.
Verdicts differ because STL cannot express BAUD-v1's
$K$-of-$M$-in-$W$ predicate exactly; the difference is informative,
not a bug.}
\label{tab:rtamt-bench}
\end{table}

\textbf{Honest reading.} (1) The verdicts differ because RTAMT's
STL approximation is strictly weaker than BAUD-v1; STL cannot count
$K$-of-$M$-in-$W$. RTAMT therefore flags the approximation while
Ghost holds on the exact property. (2) Wall-clock times measure the
verification step only; signal extraction adds $\sim$20\,ms of MCAP
parsing to RTAMT's end-to-end cost, putting it in the same order of
magnitude as Ghost. (3) The tools are complementary: RTAMT for
declarative STL over arbitrary signals, Ghost for content-addressed
CLI verification of a specific supervisor with hand-stated
predicates.

\subsection{Determinism across replicates and machines}
Within a single machine, replicate runs produce byte-identical
MCAPs. \textbf{Cross-machine byte-equality of the MCAP} is enforced
by a CI job: the reference smoke runs on a \{\texttt{ubuntu-latest},
\texttt{windows-latest}\} matrix, each runner publishes the SHA-256
of its MCAP and of its canonicalised property-report JSON (sorted
keys, \texttt{mcap\_path} dropped to absorb platform path-separator
differences), and the aggregator step \texttt{diff}s the two files.
Any disagreement fails the build.

% =======================================================================
\section{Limitations and threats to validity}

\begin{itemize}
  \item \textbf{Sim, not hardware.} MCAPs come from a simulated
        overconfidence trap, not real flight logs. Real-world
        claims require a HAL backend and hardware campaign,
        deferred to a later phase.
  \item \textbf{Reference policies only.} The TLA+ proofs and
        property semantics target specific reference policies.
        Non-reference policies would need their own ADRs, verifier
        specialisations, and TLA+ specs.
  \item \textbf{Bounded TLC.} The TLA+ proofs are exhaustive over
        finite state spaces at small constants; behaviour at
        production-scale constants rests on the property tests.
  \item \textbf{Python\,$\leftrightarrow$\,TLA+ bridge by
        inspection.} A future divergence could silently weaken the
        claim. Mitigation: review and re-run TLC on every change
        to the reference policies.
  \item \textbf{Statistical FPB out of scope.} FPB-v1 is
        observational; a statistical FPB-v2 with Monte Carlo
        bounds is a candidate future ADR.
\end{itemize}

% =======================================================================
\section{Future work}

\begin{itemize}
  \item TLA+ spec for FPB-v1 closing the formal-verification
        coverage of the full property set.
  \item TLAPS proof of the unbounded version of the partition
        theorem and of Theorem~\ref{thm:rlb}.
  \item HAL backend campaign lifting the reproducibility surface
        from simulation to flight logs.
  \item Conformance suite populating the \texttt{conformance}
        pytest marker with the HAL contract.
\end{itemize}

% =======================================================================
\section{Conclusion}
Project Ghost is a \textbf{citation pattern for safety claims}. Its
load-bearing contribution is the proposition that a safety claim
should be issued together with everything a third party needs to
reject it: a binding written statement (the ADR), a content-addressed
log (the MCAP), a pure-function verifier exposed as a CLI, property
tests, mechanical checks on the underlying invariants (TLA+/TLC),
and a signed distribution channel (OIDC PyPI). We do not claim to
have discovered the individual ingredients --- each is standard
practice --- but we have not located a tool that ships them as one
coherent unit reachable from a single shell command in the literature
we surveyed, and we treat that combination as the proper unit of
analysis.

The reference instantiation on a five-property autonomy supervisor
exists as evidence that the pattern is workable, not as the
contribution itself. The properties are deliberately simple; the
recovery latency bound is elementary in hindsight; the partition
theorem is structurally obvious. What is not obvious --- what
motivated the paper --- is that the \emph{assembly} makes safety
claims operationally falsifiable in a way that prose-and-video
assertions are not.

The artifact is re-runnable from
\texttt{pip install project-ghost==0.2.2}; the contribution stands
or falls on whether the resulting verdict means what the ADR says
it means, and on whether a third party can issue that verdict
without our help. Both, we believe, are now true.

% =======================================================================
\bibliographystyle{plainurl}
\bibliography{refs}

\end{document}
